\subsection{Behavior cloning and response masking}

SFT presents a condition and an expert continuation, then increases the continuation's
probability. A chat template first serializes system, user, assistant, and tool turns into
one causal token sequence. Prompt and tool tokens remain visible as context, but a
response-only mask normally assigns direct loss only to assistant tokens. If $m_{b,t}$
marks those targets and $M=\sum_{b,t}m_{b,t}$, the loss is
\begin{equation}
  \mathcal{L}_{\mathrm{SFT}}
  =-\frac{1}{M}\sum_{b,t}m_{b,t}
    \log p_\theta(y_{b,t+1}\mid y_{b,\leq t}).
  \label{eq:sft-loss}
\end{equation}
This is ordinary next-token cross-entropy with target selection, not a new model
architecture. $M$ counts supervised tokens rather than padded sequence length. Masking
does not avoid the forward computation for prompt context; it only decides which predicted
positions contribute directly to the scalar loss.

Token averaging gives a long assistant response more total influence than a short one.
Per-example averaging, task-balanced sampling, or explicit weights encode different
choices about whether tokens, demonstrations, or tasks should be equal units. Packing and
length-grouped batching improve utilization but do not change this statistical weighting.
